\documentclass[11pt]{book}
\usepackage{fontspec}
\setmainfont{Latin Modern Roman}
\setsansfont{Latin Modern Sans}
\usepackage{amsmath,amssymb,mathtools}
\usepackage{unicode-math}
\setmathfont{Latin Modern Math}
\usepackage{xurl}
\usepackage[hidelinks]{hyperref}
\newcommand{\sref}[2]{\hyperlink{source:#1:#2}{[S#2]}}
\newcommand{\cataloguescope}[1]{\paragraph{Mathematical question.}#1}
\begin{document}
% UnsolvedMath ID 2304028; source catalogue code AMR-022-4028
\setcounter{section}{2304027}
\section{Research Problems in Function Theory — Problem 4.28}\label{2304028}
{\small\sffamily UnsolvedMath ID 2304028 · Analysis}
\subsection{Definitions and mathematical statement}

\paragraph{Shared notation.}$\mathbb N=\{1,2,\ldots\}$, and $\mathbb Z,\mathbb Q,\mathbb R,\mathbb C$ denote the integers, rationals, reals and complex numbers. Any other conventions are specified locally.


Let $P\in\mathbb R[z]$ have degree $d\ge2$, with no assumption that all of its zeros are real. Put $H(z)=P(z)^2+P'(z)$, a polynomial of degree $2d$. Let $N_{\mathrm{nr}}(H)$ count its zeros in $\mathbb C\setminus\mathbb R$ with their algebraic multiplicities.
\paragraph{Statement recorded in the supplied source.}
Must $H$ have non-real zeros? More strongly, is
\[N_{\mathrm{nr}}\bigl(P^2+P'\bigr)\ge d-1\]
always true? Include arbitrary real-coefficient $P$, rather than only the source's earlier special case where all zeros of $P$ are real. \sref{2304028}{2}
\subsection{Short English statement}
Does the square of every nonlinear real polynomial plus its derivative have non-real zeros, with at least degree minus one counted with multiplicity?
\subsection{Sources}
\begin{description}
\item[\hypertarget{source:2304028:1}{S1}] ULAM.AI and the UnsolvedMath Contributors, \emph{UnsolvedMath: A Curated Collection of Open Mathematics Problems}, record 2304028 (AMR-022-4028). CC BY 4.0. \url{https://huggingface.co/datasets/ulamai/UnsolvedMath}.
\item[\hypertarget{source:2304028:2}{S2}] Walter K. Hayman and Edward F. Lingham, \emph{Research Problems in Function Theory}, arXiv:1809.07200, 2018. Problem 4.28, its complete Update 4.28, and the relevant chapter notation. \url{https://arxiv.org/pdf/1809.07200}.
\end{description}
\label{end:2304028}
\end{document}
